% ------------------------------------------------------------
% Chapter 1: Introduction
%
% Order is Motivation, Aims, Scope, Contributions, Terminology, Structure.
% Scope sits before Contributions so that what the study does not claim is
% fixed before what it claims is stated.
%
% Regulation is named here and argued in Section 2.2.3. Provider policy is not
% described here at all. The abstract states the findings once; Section 1.4
% states them once more with the contribution attached, and nowhere else in the
% front matter.
% ------------------------------------------------------------
\chapter{Introduction}
\label{ch:introduction}

Generative Artificial Intelligence (AI) chatbots are general-purpose tools for
writing, learning and social interaction, and children and adolescents are among
their users. The features that make them engaging, such as personalised
responses, retained conversational context and reactions to sensitive
disclosures, also raise safety concerns. For a child, an exchange appropriate for
an adult can carry age-inappropriate advice, misinformation, privacy risks, or
content that is difficult to evaluate critically at their developmental stage
\cite{figueiraActionsSpeakLouder2026}.

\section{Motivation}
\label{sec:intro_motivation}

Age gates are one layer of protection and not the only one. Prior work shows that
chatbots can infer a user's age from explicit disclosures and conversational
cues, yet may take no protective action once a child is identified
\cite{figueiraActionsSpeakLouder2026}. If a child reaches the model despite the
gate, the safety of the interaction depends on the reply itself.

Existing child-safety benchmarks evaluate harmful content and age-appropriate
language separately, and usually within single-turn exchanges. None varies the
user while holding the request fixed, so whether a reply adapts to an age the
model can already detect is unmeasured.

This thesis studies \emph{age conditioning}: the systematic change in a model's
response as a function of the age the user states, measured on safety behaviour
and on reading level. The term is behavioural throughout. It says that the reply
changes with the stated age; it does not say that the model inferred the age,
represented it, or applied any rule.

\section{Aims and Objectives}
\label{sec:intro_aims}

This thesis builds a controlled benchmark that isolates age conditioning and uses
it to test whether LLM chatbots adapt their safety behaviour and language to a
user's apparent age. The objectives are to:

\begin{enumerate}
\item construct a benchmark in which the age signal varies while the request
content is held constant;
\item score safety behaviour and linguistic accessibility on the same reply;
\item compare explicit age disclosure against implicit age cues;
\item measure whether age-conditioned behaviour survives multi-turn attack; and
\item interpret the results against provider policy and the regulatory framework
for online child safety.
\end{enumerate}

Two experiments carry them. Experiment~1, Adaptation, holds each request fixed
and varies only the age signal, covering the first three objectives.
Experiment~2, Retention, continues a stratified subset of those exchanges under
attack, covering the fourth. Separating the two keeps the single-turn comparison
matched. The fifth objective frames the discussion rather than forming a further
experiment.

\noindent
Four research questions follow:

\noindent\textbf{RQ1:}
Does a model refuse differently when the user states a child's age rather than an
adult's?

\noindent\textbf{RQ2:}
Does it write at a different reading level?

\noindent\textbf{RQ3:}
Is the change smaller when the age is implied than when it is stated?

\noindent\textbf{RQ4:}
Does the change survive an adversarial conversation?

\section{Scope}
\label{sec:intro_scope}

The study examines measurable properties of model outputs and nothing else. It
uses no data from real children, does not study user perception, and does not
identify the internal cause of any behaviour. The age is supplied rather than
elicited, so nothing here establishes that a model inferred it. Findings indicate
whether chatbots show consistent age-sensitive behaviour and where it is uneven,
not why a given model behaves as it does.

\section{Contributions}
\label{sec:intro_contributions}

The principal contribution is a controlled evaluation framework and a benchmark
of matched, age-conditioned prompts. It scores safety behaviour and linguistic
accessibility on the same reply, two properties usually assessed in isolation,
and records over-restriction alongside unsafe permissiveness so that both failure
directions are visible.

Applied to six commercial models over 46,800 requests, it establishes that age
conditioning is real and uneven. Refusal on age-restricted requests is 45.5
percentage points higher for a stated minor age than a stated adult age, and text
for a stated minor sits 1.77 grade levels lower. The two have different shapes:
almost half of the refusal movement across fourteen years of age falls in the
single step from seventeen to eighteen, where the same step carries between 4 and
11 per cent of the reading-level movement. A model can therefore be calibrated to
the statutory line and not to the developmental range beneath it. Conditioning is
also about three times weaker when the age is implied than when it is stated.

Two secondary contributions follow from the instrument. Crossing the decision
with the delivered content separates a permissive model from one that agrees and
then supplies nothing, which a refusal rate reports identically. And decomposing
the reading-level change shows it is carried by word choice rather than sentence
length.

None of the individual ingredients is new. Prior work has held prompts constant
while varying age \cite{ozcanEvaluatingChildAppropriateLanguage2026}, evaluated
demographic bias \cite{benkiraneHowCanWe2025}, contrasted explicit and implicit
cues \cite{figueiraActionsSpeakLouder2026, giorgiExplicitImplicitLarge2024},
characterised multi-turn degradation
\cite{caoSafeDialBenchFineGrainedSafety2026, liStateDependentSafetyFailures2026},
measured readability \cite{valentiniAutomaticGenerationSimplification2023,
oshikaSimplifyingTranslationsChildren2024}, and evaluated mixed sets of
open-source and closed-source models
\cite{jiaoSafeChildLLMDevelopmentalBenchmark2025}. The contribution is their
controlled integration, so that a difference in the reply is attributable to the
age signal rather than to its wording.

\section{Terminology}
\label{sec:intro_terminology}

Three terms are used for the age bands, and only in these senses. A \emph{child}
is a user aged up to 12, an \emph{adolescent} a user aged 13 to 17, and an
\emph{adult} a user aged 18 or over. Where a statement applies to both of the
first two groups they are named together. Models are \emph{closed-source} where
the weights are not released and the model is reached through a provider
endpoint, and \emph{open-source} where the weights are published and decoding can
be controlled directly.

\section{Thesis Structure}
\label{sec:intro_structure}

Chapter~\ref{ch:background} reviews related work on child safety in large
language models, age conditioning, age-appropriate language generation,
multi-turn safety and regulation. Chapter~\ref{ch:methods} describes the
benchmark, the model panel, the collection protocol and the evaluation measures,
and specifies the two experiments. Chapter~\ref{ch:results} reports the results
and Chapter~\ref{ch:discussion} interprets them. Chapter~\ref{ch:conclusion}
concludes and Chapter~\ref{ch:future} sets out future work.
